\addcontentsline{toc}{subsubsection}{Solution 11. The Harmonic--Product Squeeze}

\begin{dlsummary}[title={\bfseries Solution 11. The Harmonic--Product Squeeze},breakable]

\begin{proof}
Define
\[
g(t)=\frac1{f(1/t)}\qquad(t>0).
\]
Taking reciprocals in the first inequality and substituting
$x=1/u$ and $y=1/v$ gives
\[
g(u+v)\geq g(u)+g(v). \tag{1}
\]
The second inequality, applied to $1/u$ and $1/v$ and then reciprocated,
gives
\[
g(uv)\leq g(u)g(v). \tag{2}
\]
Thus $g$ is positive, superadditive, and submultiplicative.

\medskip
\noindent
\textbf{Step 1. Obtain a lower bound for $g$.}

Taking $x=y=1$ in (2) gives $g(1)\leq g(1)^2$, so $g(1)\geq1$.
Also, (1) and positivity show that $g$ is strictly increasing: if
$y>x>0$, then
\[
g(y)=g(x+(y-x))\geq g(x)+g(y-x)>g(x).
\]
Fix $t>1$ and a positive integer $n$. Submultiplicativity and
superadditivity give
\[
g(t)^n\geq g(t^n)\geq\lfloor t^n\rfloor g(1)
\geq\lfloor t^n\rfloor.
\]
The middle inequality follows by writing $t^n$ as the sum of
$\lfloor t^n\rfloor$ copies of $1$ and, if needed, a positive
remainder. Taking $n$th roots and letting $n\to\infty$ yields
\[
g(t)\geq t\qquad(t>1). \tag{3}
\]

Let $a=1/c>1$. By the definition of $g$,
\[
g(a)=\frac1{f(c)},\qquad \frac1{g(a)}=f(c),
\]
and therefore
\[
g(g(a))=\frac1{f(1/g(a))}
       =\frac1{f(f(c))}\leq\frac1c=a.
\]
On the other hand, (3) applied first to $a$ and then to $g(a)>1$ gives
\[
g(g(a))\geq g(a)\geq a.
\]
Therefore,
\[
g(a)=a. \tag{4}
\]

\medskip
\noindent
\textbf{Step 2. Establish a global upper bound.}

By (2) and (4), $g(a^n)\leq g(a)^n=a^n$ for every positive integer
$n$. Fix $x>0$ and, for sufficiently large $n$, put
\[
m_n=\left\lfloor\frac{a^n}{x}\right\rfloor\geq1.
\]
Since $m_nx\leq a^n$, monotonicity and repeated superadditivity give
\[
a^n\geq g(a^n)\geq g(m_nx)\geq m_ng(x).
\]
Consequently,
\[
g(x)\leq\frac{a^n}{\lfloor a^n/x\rfloor}\longrightarrow x,
\]
so $g(x)\leq x$ for all $x>0$.

\medskip
\noindent
\textbf{Step 3. Force equality.}

The upper bound gives $g(1)\leq1$, and hence $g(1)=1$. By (2),
\[
1=g(1)=g(x\cdot1/x)\leq g(x)g(1/x)\leq\frac{g(x)}x.
\]
Thus $g(x)\geq x$. Together with the upper bound, this proves
$g(x)=x$ for all $x>0$. Returning to the definition of $g$,
\[
\frac1{f(1/x)}=x,
\]
so $f(x)=x$ for every $x>0$.

\medskip
\noindent
\textbf{Verification.}

For $f(x)=x$, both inequalities hold with equality, and
$f(f(c))=c$. Hence the identity function is indeed the unique solution.
\end{proof}

\end{dlsummary}